\section{The Big Comeback: Deep Learning Sees the World}

\subsection{The ImageNet Moment (2012)}

By 2010, three things had quietly converged: larger datasets, faster GPUs, and better training techniques.
The spark that lit the fire came in 2012.

Every year, researchers competed in the \textbf{ImageNet Large Scale Visual Recognition Challenge} \cite{russakovsky2015imagenet} --- a contest to correctly classify one million images into one thousand categories.
In 2012, a team from the University of Toronto, led by Geoffrey Hinton and his students Alex Krizhevsky and Ilya Sutskever, entered a deep neural network called \textbf{AlexNet} \cite{krizhevsky2012imagenet}.

The result was staggering.
AlexNet cut the error rate nearly in half compared to the previous year's best method.
Nothing like this had ever happened in the competition's history.

Overnight, the machine learning world shifted.
Deep learning was no longer a niche academic interest --- it was the new state of the art.

\subsection{What Made It Work: Convolutional Neural Networks}

AlexNet used a special architecture called a \textbf{Convolutional Neural Network} (CNN).
The core idea behind CNNs is that images have \emph{spatial structure}: nearby pixels are related, and patterns like edges and shapes appear in many places across an image.

A standard neural network treats every pixel independently.
A CNN, by contrast, uses small filters that slide across the image looking for local patterns --- first detecting edges, then shapes, then objects.
Each layer builds on the last, from simple to complex.

\begin{tcolorbox}[colback=blue!5!white, colframe=blue!40!black, title=How a CNN Sees a Cat]
\textbf{Layer 1:} Detects edges --- horizontal lines, vertical lines, diagonals. \\
\textbf{Layer 2:} Combines edges into curves, corners, and textures. \\
\textbf{Layer 3:} Recognizes shapes --- eyes, ears, whiskers. \\
\textbf{Layer 4:} Assembles shapes into the concept of a ``cat.''
\end{tcolorbox}

CNNs were not entirely new --- Yann LeCun had pioneered them in the late 1980s for reading handwritten digits on cheques \cite{lecun1989backpropagation}.
But AlexNet showed that with enough data and GPU power, they could tackle the full complexity of natural images.

\subsection{The Gold Rush}

The years following 2012 saw an explosion of investment and research.
Google, Facebook, Microsoft, and Baidu all built dedicated AI research labs.
Deep learning models rapidly surpassed human-level performance on image recognition benchmarks.
They were applied to medical imaging, autonomous vehicles, satellite imagery, and facial recognition.

The field had a tool that genuinely worked at scale.
The question that followed naturally was: \emph{what else can we apply it to?}
